
\documentclass[fleqn,head_space=14mm,foot_space=15mm,gutter=10mm,fore_edge=10mm,
    twocolumn,jafontscale=0.92]{jlreq}

\setlength{\columnseprule}{0.4pt}
\setlength{\mathindent}{2\zw}
\setlength{\fboxsep}{5pt}
\setlength{\fboxrule}{1.0pt}

%セクション見出し設定
\ModifyHeading{section}{lines=2,label_format={\S},
    format={%
        \begin{tikzpicture}
        \useasboundingbox (0,0) rectangle (0,0);
        \fill[black!50] (0,-4pt) rectangle (8pt,14pt);
        \draw[black!50,line width=1pt] (0,-4pt) -- (\linewidth,-4pt);
        \end{tikzpicture}%
    \hspace{1\zw}#1#2}}

\ModifyHeading{subsection}{font=\sffamily\large,
    label_format={\fbox{\Large{\arabic{subsection}}}}}

%ヘッダー・フッター設定
\usepackage{titleps}
    \newpagestyle{main}{
    \sethead[\textsf{【大学入試 物理】 標準演習〈問題〉 ver.01}][]
            [\textsf{\S \, \sectiontitle}]
            {\textsf{【大学入試 物理】 標準演習〈問題〉 ver.01}}{}
            {\textsf{\S \, \sectiontitle}}
    \setfoot[][][]
            {}{\thepage}{}
    }
    
    \pagestyle{main}

\usepackage[deluxe,match,jfm_yoko=jlreq,jfm_tate=jlreqv]{luatexja-preset}

\usepackage{enumitem}
\def\threedigits#1{
    \ifnum#1<100 0\fi 
    \ifnum#1<10 0\fi 
    \number#1}

\usepackage{tikz}
\usetikzlibrary{patterns,angles,quotes}

\usepackage{lua-ul}

\usepackage{amsmath}
\usepackage[lmodern]{mathcomp}

\usepackage{fancybox}

\frenchspacing

\begin{document}

\setlength{\abovedisplayskip}{0mm}
\setlength{\belowdisplayskip}{0mm}

\setlist[enumerate]{label=\textsf{（\hspace{-.3pt}\arabic*\hspace{-.3pt}）}\, ,
    topsep=2pt, itemsep=1.5pt, itemindent=1\zw, leftmargin=*}
  
\section{図形の証明}

\subsection{都立Vもぎ 2023年9月}

下の図で，四角形ABCDは，$\text{AB}<\text{BC}$の長方形である．

長方形ABCDの外部に点Eを，辺BCと線分AEが交わるようにとる．

辺BCと線分AEとの交点をFとする．

頂点Aと頂点C，頂点Cと点Eをそれぞれ結ぶ．

$\text{BF}=\text{EF}$，$\angle\text{CEF}=90\tcdegree$のとき，
$\triangle\text{ABF}\equiv\triangle\text{CEF}$であることを証明せよ．


\vspace{4pt}
\begin{center}

    \begin{tikzpicture}
        \draw[->,>=stealth,semithick] (0,0) -- (7,0) node[right] {$x$};
        \draw[->,>=stealth,semithick] (0,0) -- (0,5) node[above] {$y$};
        \draw[->,>-stealth,thick] (0,0) -- (1,1.5) node[above right=-2pt] {$v_0$};
        \draw (1,0) coordinate (A) -- (0,0) coordinate (O) -- (0.2,0.3) coordinate (B);
        \filldraw[thick,fill=black!50!white] (0,0) circle (5pt) node[below left] {O};
        \draw pic["$\theta$", draw=black, angle eccentricity=1.35, angle radius=7mm] 
            {angle=A--O--B};

    \end{tikzpicture}
\end{center}


\newpage

\section{しらん}

\newpage

\section{静電場}

\subsection{点電荷が作る電場}

真空中に$xy$平面を定め，$y$軸上の点$(0, \hspace{1.2mm} L)$に正の点電荷A（電荷$Q$）を固定し，
さらに$y$軸上の点$(0, \hspace{1.2mm} -L)$に負の点電荷B（電荷$-Q$）を固定した．
クーロンの法則の比例定数を$k_0$とし，クーロン力以外の力の影響は考えない．
点$(x, \hspace{1.2mm} y)$における電場の$x$成分を$E_x(x, \hspace{1.2mm} y)$，
$y$成分を$E_y(x, \hspace{1.2mm} y)$，
電位（無限遠基準）を$\phi(x, \hspace{1.2mm} y)$と表す．

\begin{enumerate}

    \item $E_x(L, \hspace{1.2mm} L)$，$E_y(L, \hspace{1.2mm} L)$を求めよ．

    \item $E_x(x, \hspace{1.2mm} 0)$，$E_y(x, \hspace{1.2mm} 0)$，
    $\phi(x, \hspace{1.2mm} 0)$を求めよ．

    \item $E_x(0, \hspace{1.2mm} y)$，$E_y(0, \hspace{1.2mm} y)$，
    $\phi(0, \hspace{1.2mm} y)$を求めよ．
    ただし，$y \neq \pm L$とする．

    \item $y$軸に沿ってなめらかに動ける正の点電荷C（質量$m$，電荷$q$）を，
    点$(0, \hspace{1.2mm} 2L)$で静かに放したところ，Cは無限遠方へ飛び去った．
    Cが無限遠に到達したときの速さ$v_\infty$を求めよ．

\end{enumerate}

\vspace{2cm}
\vfill

\subsection{平行一様電場・電位の基準}

真空中に定めた$x$軸に沿って，$x$軸の負の向きに大きさ$E$の平行一様電場が在る．
この電場中における正の点電荷（質量$m$，電荷$q$）の$x$軸に沿った運動を考える．
点電荷には，原点（$x=0$）において，点電荷に$x$軸の正の向きに初速$v_0$を与えたものとし，
クーロン力以外の力の影響は考えない．

\begin{enumerate}

    \item 点電荷が位置$x$に達したときの速度を$v$とする．
    そのときの速さ$|v|$を，$x$の関数として表せ．

    \item 位置$x$の電位$\phi(x)$を求めよ．
    ただし，$x=h$の位置を電位の基準とする．

\end{enumerate}

\vspace{2cm}
\vfill

\subsection{拡がった電荷分布の作る電場}

真空中の誘電率を$\varepsilon_0$とする．

\begin{enumerate}

    \item クーロンの法則の比例定数$k_0$を，$\varepsilon_0$を用いて表せ．
    
    \item 充分長い直線上に一様に電荷が分布しており，
    単位長さあたりの電荷（電荷の線密度）は$\lambda \mspace{5mu} (\lambda>0)$である．
    この場合，直線に充分近い点では，電気力線は直線に垂直に一様に出ていると見なせる．
    その条件下で，直線から距離$r$の点での電場の強さ$E(r)$を求めよ．

    \item 充分広い平面上に一様に電荷が分布しており，
    単位面積あたりの電荷（電荷の面密度）は$\sigma \mspace{5mu} (\sigma>0)$である．
    この場合，平面に充分近い点では，電気力線は平面に垂直に一様に出ていると見なせる．
    その条件下で，周囲での電場の強さ$E$を求めよ．

    \item 真空中に半径$a$の球殻（厚み無視）があり，一様に帯電している．
    その電荷の総量は$Q \mspace{5mu} (Q>0)$である．
    球殻の中心から$r$離れた点での電場の強さ$E(r)$を，
    $0<r<a$の領域と$r>a$の領域に分けて，それぞれ求めよ．

    \item 真空中に半径$a$の荷電球があり，一様に帯電している．
    その電荷の総量は$Q \mspace{5mu} (Q>0)$である．
    球殻の中心から$r$離れた点での電場の強さ$E(r)$を，
    $0<r<a$の領域と$r>a$の領域に分けて，それぞれ求めよ．

    \item ともに$+Q$の電荷を持つ2つの点電荷が，ある距離を隔てて固定されている．
    この周囲の電気力線の概略を図示せよ．

\end{enumerate}

\vspace{2cm}
\vfill

\subsection{電位から電場を逆算する①}

$xy$平面上の電位分布を測定し，2\,V間隔で等電位線を描いたところ，
図のようになった．
等電位線の横に添えた数値は，各等電位線上での電位の値である．

\begin{enumerate}
    
    \item 陽子（電荷$e=1.6 \times 10^{-19}$\,C）を，
    クーロン力に逆らってゆっくりと運ぶとする．
    点Pから点Fへと運ぶ際に必要な仕事$W_\text{PF}$，
    点Pから点Gへと運ぶ際に必要な仕事$W_\text{PG}$，
    点Pから点Hへと運ぶ際に必要な仕事$W_\text{PH}$を，それぞれ求めよ．

    \item 電子（電荷$-e=-1.6 \times 10^{-19}$\,C）を，
    クーロン力に逆らってゆっくりと運ぶとする．
    点Pから点Fへと運ぶ際に必要な仕事$W^*_\text{PF}$，
    点Pから点Gへと運ぶ際に必要な仕事$W^*_\text{PG}$，
    点Pから点Hへと運ぶ際に必要な仕事$W^*_\text{PH}$を，それぞれ求めよ．

    \item 点Pを通る電気力線の概略を描け．
    
    \item 点P付近に図示してある等電位線の間隔は大体0.53\,m程度である．
    このことから点Pにおける電場の強さ$E_\text{P}$の近似値を求めよ．

\end{enumerate}

\vspace{4pt}

\begin{center}

    \begin{tikzpicture}
    
        \pgfmathsetmacro{\xmax}{7}
        \pgfmathsetmacro{\ymax}{7}
        \pgfmathsetmacro{\a}{3.22}

        \draw[->,>=stealth,semithick] (-.5,0) -- (\xmax,0) node[right] {$x$};
        \draw[->,>=stealth,semithick] (0,-.5) -- (0,\ymax) node[above] {$y$};
        \node at (0,0) [below left] {O};
        
        \begin{scope}\clip(0,0)rectangle(\xmax,\ymax);
                        
            \draw[samples=500,domain=0.1:\xmax] plot(\x,\a/\x);
            \draw[samples=500,domain=0.1:\xmax] plot(\x,\a*2/\x);
            \draw[samples=500,domain=0.1:\xmax] plot(\x,\a*3/\x);
            \draw[samples=500,domain=0.1:\xmax] plot(\x,\a*4/\x);
            \draw[samples=500,domain=0.1:\xmax] plot(\x,\a*5/\x);
            \draw[samples=500,domain=0.1:\xmax] plot(\x,\a*6/\x);
            \draw[samples=500,domain=0.1:\xmax] plot(\x,\a*7/\x);
            \draw[samples=500,domain=0.1:\xmax] plot(\x,\a*8/\x);
            \draw[samples=500,domain=0.1:\xmax] plot(\x,\a*9/\x);
            \draw[samples=500,domain=0.1:\xmax] plot(\x,\a*10/\x);
            \draw[samples=500,domain=0.1:\xmax] plot(\x,\a*11/\x);
            \draw[samples=500,domain=0.1:\xmax] plot(\x,\a*12/\x);
            \draw[samples=500,domain=0.1:\xmax] plot(\x,\a*13/\x);
            \draw[samples=500,domain=0.1:\xmax] plot(\x,\a*14/\x);
         
        \end{scope}

        \pgfmathsetmacro{\xa}{2.2}
        \pgfmathsetmacro{\xb}{4.8}
        \pgfmathsetmacro{\h}{\a*2/\xb}
        \pgfmathsetmacro{\g}{\a*3/\xb}
        \pgfmathsetmacro{\f}{\a*7/\xb}
        \pgfmathsetmacro{\p}{\a*3/\xa}

        \filldraw[fill=black] (\xb,\h) circle (1.6pt) node[below] {H};
        \filldraw[fill=black] (\xb,\g) circle (1.6pt) node[above] {G};
        \filldraw[fill=black] (\xb,\f) circle (1.6pt)
            node[above=2.3mm,right=-1.4mm] {F};
        \filldraw[fill=black] (\xa,\p) circle (1.6pt)
            node[above=2.3mm,right=-1mm] {P};

        \node at (\xmax,\a/\xmax) [right] {\small{2\,V}};
        \node at (\xmax,\a*3/\xmax) [right] {\small{6\,V}};
        \node at (\xmax,\a*5/\xmax) [right] {\small{10\,V}};
        \node at (\xmax,\a*7/\xmax) [right] {\small{14\,V}};
        \node at (\xmax,\a*9/\xmax) [right] {\small{18\,V}};
        \node at (\xmax,\a*11/\xmax) [right] {\small{22\,V}};
              
    \end{tikzpicture}

\end{center}

\vspace{2cm}
\vfill

\subsection{電位から電場を逆算する②}

真空中に電場が在り，$xy$平面上での電位分布が
\[ \phi=a(x^2-y^2)\,(a=10\,\text{V/}\text{m}^2) \]
で与えられるものとする．

\begin{enumerate}
    
    \item $xy$平面上の$x>0$かつ$y>0$の領域について，
    $\phi=0$\,V，$\phi=\pm10$\,V，$\phi=\pm20$\,Vの
    各々に対応する等電位線の概形を図示せよ．
    また，同じ図上に点$(1\,\text{m}, \hspace{1.2mm}1\,\text{m})$
    を通る電気力線の概形を描け．

    \item 点$(1\,\text{m}, \hspace{1.2mm}0\,\text{m})$での
    電場の強さと向きを求めよ．

    \item $y$軸に沿ってなめらかに動ける正の点電荷P（質量$m$，電荷$q$）に対し，
    点$(0, \hspace{1.2mm}-h)\,(h>0)$で初速度（$y$軸の正の向きに大きさ$v_0$）を与える．
    Pが$y$軸正方向の無限遠へ飛び去るための$v_0$の条件を求めよ
    （この設問のみ文字式で答えよ）．

\end{enumerate}

\vspace{2cm}
\vfill

\subsection{点電荷が作る電場}

真空中に$xy$平面を定め，$y$軸上の点$(0, \hspace{1.2mm} L)$に
電荷$Q\,(Q>0)$の点電荷Aを固定し，さらに$y$軸上の点$(0, \hspace{1.2mm} -L)$に
同じく電荷$Q$の点電荷Bを固定した．クーロンの法則の比例定数を$k_0$とし，
クーロン力以外の力の影響は考えない．
点$(x, \hspace{1.2mm} y)$における電場の$x$成分を$E_x(x, \hspace{1.2mm} y)$，
$y$成分を$E_y(x, \hspace{1.2mm} y)$，
電位（無限遠基準）を$\phi(x, \hspace{1.2mm} y)$と表す．

\begin{enumerate}

    \item $E_x(x, \hspace{1.2mm} 0)$，$E_y(x, \hspace{1.2mm} 0)$，
    $\phi(x, \hspace{1.2mm} 0)$を求めよ．

    \item 正の点電荷P（質量$m$，電荷$q$）を，無限遠から
    点$(\sqrt{3}L, \hspace{1.2mm} 0)$までゆっくりと運ぶ．
    この際に必要な仕事$W$を求めよ．

    \item Pが$x$軸上をなめらかに動けるものとする．
    Pに点$(\sqrt{3}L, \hspace{1.2mm} 0)$で，$x$軸負方向に初速$v_0$を与える場合，
    Pが$x$軸負方向の無限遠へ飛び去るための$v_0$の条件を求めよ．

    \item 負の点電荷N（質量$m$，電荷$-q$）を，無限遠から
    点$(\sqrt{3}L, \hspace{1.2mm} 0)$までゆっくりと運ぶ．
    この際に必要な仕事$W^*$を求めよ．

    \item Nが$x$軸上をなめらかに動けるものとする．
    Nに点$(\sqrt{3}L, \hspace{1.2mm} 0)$で，$x$軸正方向に初速$v_0$を与える場合，
    Nが$x$軸正方向の無限遠へ飛び去るための$v_0$の条件を求めよ．

    \item Nを原点付近の$x$軸上の点で静かに放したところ，Nの運動は単振動と見なせるものとなった．
    その周期を求めよ．ただし，微小量の2次の項は無視せよ．

\end{enumerate}

\vspace{2cm}
\vfill

\end{document}